%---------------------------------------------------------------------
\chapter{Requirements and System Analysis}
\label{ch:req}
%---------------------------------------------------------------------

\section{Stakeholders and Users}
The primary user is the \textbf{compliance officer} of a producer firm, who owns the reconciliation and
who must defend the producer's position. This user needs to know which discrepancy deserves attention
first and why, with every figure tied to a record and every rule tied to a dated provision.
The second user is the \textbf{reviewer}, who may be the same person or a senior colleague, and who must
decide the report. The reviewer needs to see exactly what evidence and which calculation versions the
report rests on, and needs a decision that is bound to that version, so that a later change to the report
cannot be approved by an earlier decision.

Secondary stakeholders are the \textbf{recyclers} whose certificates are examined, the \textbf{authorised
signatory} of the producer who files the return, and the \textbf{auditors} or \textbf{regulators} who may
later ask how a position was reached. None of these is a direct user of the prototype, but their needs
explain the requirement for a complete audit trail and for a clear statement of what the system did not
check.

\section{Functional Requirements}
\label{sec:fr}
Each requirement is stated so that it can be tested, and each is implemented in the repository.
\begin{description}[itemsep=2pt,font=\normalfont\bfseries,leftmargin=1.2cm,labelwidth=1cm]
\item[FR1] The system shall accept a case dossier described by a manifest of documents and records and
           shall record an intake outcome, including failures, for every file.
\item[FR2] The system shall store each original file unchanged, content-addressed by its SHA-256 digest,
           and shall never modify an original; corrections create linked revisions.
\item[FR3] The system shall extract granular observations from PDF, CSV and JSON files, each with an exact
           locator (page, row or JSON pointer), and shall distinguish missing, unknown, not applicable and zero.
\item[FR4] The system shall normalise units and entity aliases and shall refuse to add quantities of
           different dimensions or different metals.
\item[FR5] The system shall freeze a case snapshot and a corpus snapshot with digests, and a run shall pin
           the exact snapshots, calculation profile, index and prompt it used.
\item[FR6] The system shall compute, for each obligation bucket, the required quantity, the supported
           quantity, the signed gap and the pending quantity, using exact decimal arithmetic.
\item[FR7] The system shall retrieve regulatory provisions from a dated corpus, filter them for
           temporal applicability and authority, and report the source, version and locator of every
           provision it relies on.
\item[FR8] The system shall run a bounded investigation in which a model proposes a check, a tool executes
           it, and the model assesses the observed result, until it resolves, abstains or stops.
\item[FR9] The system shall accept a model-proposed finding only if it passes deterministic checks of its
           citations, numbers and wording, and shall retain every rejected proposal with its reason.
\item[FR10] The system shall build a structured report, verify its numbers and citations, and block the
           report if verification fails after one repair.
\item[FR11] The system shall present the report to an authenticated reviewer and shall bind the reviewer's
           decision to the exact report and input versions.
\item[FR12] The system shall write an append-only audit trail of every plan, model call, tool call,
           observation and decision, and shall export a verifiable bundle from which a run can be replayed.
\item[FR13] The system shall provide an evaluation harness that runs variants on dossiers, scores them
           against an independent reference in a separate process, and produces tables with explicit
           denominators.
\end{description}

\section{Non-Functional Requirements}
\begin{description}[itemsep=2pt,font=\normalfont\bfseries,leftmargin=1.2cm,labelwidth=1cm]
\item[NFR1] \textbf{Exactness.} No compliance quantity shall be computed with binary floating point, and no
            quantity shall be produced by the model.
\item[NFR2] \textbf{Auditability.} The intent of every external call shall be recorded before it is made, and
            every attempt shall be retained, including failed and uncertain ones.
\item[NFR3] \textbf{Safety.} The model shall have no write access, no shell, no arbitrary query and no
            ability to approve evidence or impersonate a reviewer. Dossier and retrieved text shall be
            treated as untrusted data.
\item[NFR4] \textbf{Locality.} Dossier data and model calls shall stay on the machine; the model endpoint
            shall be loopback only.
\item[NFR5] \textbf{Boundedness.} Each run shall have fixed limits on model calls, tokens, tool calls per
            finding, reformulations, retries and wall time, enforced before a call is dispatched.
\item[NFR6] \textbf{Reproducibility.} Runs shall pin their inputs, and the portable parts of the system
            shall be testable on Windows, Linux and macOS without the model.
\item[NFR7] \textbf{Honesty.} Abstention, qualification and error outcomes shall be preserved and reported,
            and undefined metrics shall be reported as not applicable and not as zero or one.
\end{description}

\section{Data Sources}
\label{sec:data}
Three kinds of data are used.

\textbf{Synthetic dossiers.} The evaluation uses 48 fictional dossiers generated from a typed world model
with exact decimal quantities (Section~\ref{sec:dataset}). Each dossier holds rendered PDFs (invoices and
memos), CSV ledgers, and JSON records for entities, obligations, certificates, allocations, balances and
registrations. All names and identifiers are fictional and every document is watermarked as synthetic.

\textbf{Regulatory corpus.} Retrieval uses two source documents: the E-Waste (Management) Rules, 2022
(38 physical pages) and the CPCB frequently-asked-questions document of 23 January 2024 (13 pages). The
fixture corpus retains the pages for Rule~6(4) and Rule~15(1)--(3) of the Rules and the page of the FAQ on
recycler capability. Both sources have a recorded rights state of \emph{review required}, and the release
plan is to publish their identifiers and digests rather than the files. The small size of this corpus limits
any claim about retrieval quality.

\textbf{Calculation profile.} A profile declares the allowed metals, the target basis, the unsupported
features and the state of the purchase ceiling of Rule~15(1), which stays disabled until its applicability
is reviewed. An ordinary supplied target remains usable even when a refurbishment offset is unsupported.

\section{Hardware and Software Requirements}
\begin{table}[H]
\centering
\caption{Environment used for the development and the real-model runs}
\label{tab:env}
\small
\begin{tabularx}{\linewidth}{@{}l X@{}}
\toprule
\textbf{Item} & \textbf{Value} \\ \midrule
Machine & Apple silicon laptop, 16\,GB unified memory \\
Model & \code{mlx-community/Qwen3.5-9B-4bit}, revision \code{8b2b98c}, MLX-LM 0.32.2, loopback port 8766 \\
Generation speed & about 8.7 tokens per second (measured on this machine) \\
Backend language & Python 3.11 or newer, managed with \code{uv} \\
Main libraries & LangGraph 1.x with the SQLite checkpointer, FastAPI, Pydantic 2, pdfplumber, Typer, Uvicorn \\
Storage & SQLite (application database, audit kernel, FTS5 index) and a content-addressed blob store \\
Interface & React 19 with TypeScript and Vite, served against the HTTP API \\
Quality tools & pytest (1230 tests collected), Ruff for linting and formatting \\
\bottomrule
\end{tabularx}
\end{table}

\section{Feasibility}
\textbf{Technical.} All components run on one laptop. The model is the limiting component: at about
8.7 tokens per second, a run that makes several calls takes one to three minutes. This is acceptable for
an advisory working paper, and it is the reason the evaluation is sized in tens of dossiers and not hundreds.

\textbf{Data.} Real producer data is not available and would raise confidentiality issues. Synthetic
dossiers with an independent reference allow every claim to be scored. The cost is that synthetic data
cannot show behaviour on real records, which Chapter~\ref{ch:limits} states.

\textbf{Operational and financial.} The model is local, so the provider cost of every run is zero. The
compute and energy used are not measured, and the report records them as not measured and not as zero.

\textbf{Schedule.} The reviews fall on 16 October (charter), 24 October (design), 31 October and 7 November
2026, with final artifacts due on 12 November 2026. The held-out evaluation depends on a stable prompt and
model setup and is the part of the plan that is not complete in this report.
